% =============================================================================
%  s01_intro.tex — Sección 1: Introducción / Planteamiento del problema
%  Sustituye \lipsum con tu contenido real.
% =============================================================================

\section{Introducción}


% ── Frame: Solo texto — mensaje clave, cita o puente entre secciones ─────────
%    USO: para declaraciones fuertes, preguntas retóricas o transiciones.
%    Copia y adapta. No lleva título de frame ni elementos gráficos.
\begin{frame}[plain]
  \begin{tikzpicture}[remember picture, overlay]
    % Banda lateral de color
    \fill[ColorPrincipal]
      (current page.north west) rectangle
      ([xshift=0.35cm] current page.south west);
    % Línea dorada en el borde de la banda
    \fill[ColorAcento]
      ([xshift=0.32cm] current page.north west) rectangle
      ([xshift=0.35cm] current page.south west);
  \end{tikzpicture}
  \hspace{0.8cm}%
  \begin{minipage}{0.88\textwidth}
    \vspace{2.2cm}
    {\color{ColorAcento}\rule{0.5\textwidth}{1.2pt}}\\[0.5cm]
    {\large\bfseries\color{ColorPrincipal}
      Pregunta de Investigación}\\[0.6cm]
    {\normalsize\color{black}

    ¿En qué medida la optimización topológica de la red de transporte de la Ciudad de
    México y la Zona Metropolitana del Valle de México, evaluada mediante el motor analítico
    VFTModel \citep{garibaldi2023vft}, demuestra que la implementación de corredores anillares 
    reduce la vulnerabilidad funcional y mejora la eficiencia global del sistema en comparación
    con el modelo radial vigente?}\\[0.5cm]
    {\color{ColorAcento}\rule{0.3\textwidth}{0.8pt}}
    \vfill
  \end{minipage}
\end{frame}

% ── Frame: Objetivos ─────────────────────────────────────────────────────────
\begin{frame}{Objetivos de la Presentación}
  \begin{block}{Objetivo General del Taller}
    En que medida los temas vistos durante el semestre 2026-2, benefician y nutren la investigación
    sobre transportes anillares bajo 3 ejes: \textbf{Social, Económico y Social} en el entorno de una
    Megalópolis como lo es la Ciudad de México y su área Metropolitana.
  \end{block}
  \vspace{0.4cm}
  \begin{columns}[T]
    \begin{column}{0.32\textwidth}
      \centering
      {\color{ColorPrincipal}\Large\textbf{Social}}\\[0.1cm]
      {\small Contexto Social}
    \end{column}
    \begin{column}{0.32\textwidth}
      \centering
      {\color{ColorPrincipal}\Large\textbf{Económico}}\\[0.1cm]
      {\small Contexto Económico}
    \end{column}
    \begin{column}{0.32\textwidth}
      \centering
      {\color{ColorPrincipal}\Large\textbf{Geográfico}}\\[0.1cm]
      {\small Contexto Geográfico}
    \end{column}
  \end{columns}
\end{frame}
